\subsubsection{Comparison with the Acting Master Plan of Gatchina}
\label{sec:comparison-genplan}
 
The demonstration with Gatchina was discussed above as an illustration of the FP2MP coordination mechanism on an open-ended task in which no single correct solution exists. The same artefact admits a different reading: it can be compared against the acting master plan of the same city---the territorial-planning document that fixes facility locations, functional zones, and built-form parameters. Such a comparison answers the question of whether the FP2MP proposal extends beyond an analytical illustration and whether it aligns with the real planning context. This question matters because the weakly formalized tasks for which FP2MP is positioned are, in actual planning practice, resolved not through arbitrary re-evaluation but within an existing institutional frame: an approved master plan that has passed public hearings and a regulatory compliance check. We performed this comparison and record both its results and its limitations below.
 
\subsubsection{Comparison Frame}
\label{sec:comparison-frame}
 
The comparison is conducted over two documents:
 
\begin{itemize}
  \item the acting Master Plan of the municipal formation ``City of Gatchina'' (revised edition)---an approved territorial-planning document with a horizon of 2040; a design population of 112.0 thousand inhabitants; it includes the Regulation on Territorial Planning and the Norms of Urban-Planning Design;
  \item the FP2MP proposal---the artefact produced by the demonstration run: 351 blocks, 83\,320 inhabitants in the local model.
\end{itemize}
 
The comparison was performed on real proposals extracted from the sources rather than on term frequencies. In total, 109 content-bearing statements were aligned: 72 from the master plan (146 facilities of the first part and 34 facilities of part~2) and 37 from FP2MP (three implementation waves, five territory types, eight priority residential blocks, five service centres, the 196-node transport anchor, and the transport action block).
 
We focus on the master plan rather than on the Strategy of Socio-Economic Development, because FP2MP operates on a local city model and spatial prioritisation---that is, at the same level as the master plan. The Strategy sets the political-economic frame (mission, budget, demographic scenario, EVENT-area programme) and its direct comparison with FP2MP blurs the positioning of the method: 30 of the 47 content-bearing statements of the Strategy have no direct relation to spatial prioritisation.
 
\subsubsection{Comparison by Topic}
\label{sec:comparison-topics}
 
The documents diverge already at the population level relative to which they formulate their proposals: the master plan is 112.0 thousand by 2040, whereas FP2MP is 83.3 thousand in the local model. This gap is not a technical error and cannot be resolved by choosing the ``right'' number. These are two different operational bases.
 
FP2MP operates with 83.3 thousand not because it ignores official data, but because the BlocksNet block model contains exactly this number. The master plan selects the maximum of the available trajectories and serves its infrastructure frame accordingly. If FP2MP is to move from a local model to a real planning instrument, it must select a demographic scenario and justify the choice. Without this, its eight priority residential blocks (block ids 331, 108, 289, 87, 89, 9, 18, 2---see Section~5.2.4) serve 35\,991 inhabitants in the model, but do not answer the question of where the additional 28\,680 inhabitants foreseen by the master plan for 2040 are to be housed.
 
The master plan allocates 683.3~ha of housing: 60.7~ha of multi-storey (12 storeys, 40\% built-up, 6\,100\,m\textsuperscript{2}/ha gross), 135.5~ha of mid-rise, 59.1~ha of low-rise, 428~ha of individual housing. This is the resource frame for 112.0 thousand inhabitants. FP2MP, in its ``Options I weighed'' section, explicitly rejects mass peripheral expansion: the peripheral block~285 with an area of 1\,678\,249\,m\textsuperscript{2} is not treated as a new-housing zone even in the third wave (7--12 years).
 
What we observe here is not a contradiction but a unit mismatch. The master plan speaks in hectares and built-form parameters, FP2MP in blocks and priorities. Alignment requires translation: the eight priority residential blocks of FP2MP occupy a certain share of one of the master plan housing zones. Until such a cross-check is performed, claims of ``match'' or ``conflict'' are premature. However, the very existence of FP2MP as a ``rejection of expansion'' is structurally incompatible with a master plan in which expansion is already locked in. This is a mismatch of objectives, not an error of FP2MP as such. We attribute it to a difference in operational models: the master plan normatively fixes the future, FP2MP diagnoses the present.
 
The master plan contains 83 specific facilities with parameters. Representative examples: replacement of the ``Elteza'' boiler house with a new BMK of 4.0~Gcal/h at ul. Matveeva~48; construction of a 24.5~Gcal/h boiler house at ul. Kirgetova~21 for residential and public-business development after the disconnection of FGUP ``PEKP''; reconstruction of boiler house No.~11 in Promzona-1 to 254.8~Gcal/h; main heating networks of 6.0~km in Zayachy Remiz and 2.9~km in V``ezd; the South-Western water intake with a water-treatment station; municipal wastewater treatment facilities of 60.0~thousand~m\textsuperscript{3}/day north of the village of Vayyalovo; 17.46~km of water-supply networks in Egerskaya Sloboda, Zagvozdka, Marienburg, and Khimozi; reconstruction of the water main with diameter increase from 700 to 1000~mm at ul. Ryseva; three 10~kV distribution points in Aerodrom, Zayachy Remiz, and at ul. Generala Knysha; high-pressure gas pipelines with GRP of 0.32--0.50~km in Roshchinsky, V``ezd, and Zayachy Remiz.
 
FP2MP does not include engineering infrastructure as an independent direction: the only substantive mention is the conditional permit for peripheral development (``pass an engineering check'').
 
This observation has substantive content: the master plan is a programme of replacement and development of engineering infrastructure, whereas FP2MP is an instrument operating on top of an existing infrastructure frame. In its current configuration FP2MP does not pose the task of network modernisation; it coordinates social and spatial policy under the condition that the infrastructure frame already exists or will be created within the master plan.
 
The master plan creates a connected network: city-wide arterials of 1.8--2.7~km, four lanes, across the microdistricts of Zayachy Remiz, Aerodrom, Promzona-2, and ul. Krasnykh Voenl\"yotov; district-level arterials in eleven microdistricts; bridges over the Belaya (Kirgetova) and Serebryanaya (Voskova) rivers, overpasses over railway tracks; the railway station ``Gatchina-Varshavskaya'' within the transport infrastructure zone; multi-storey and underground car parks of 250, 160, and 30 parking spaces in V``ezd and Centre; open-air parking of 100 and 65~spaces in Aerodrom and Centre.
 
FP2MP enumerates seven transport actions: reinforcement of bus links between residential subcentres 331, 108, 289, 87, 89, 9, 18, 2 and service centres 136, 135, 177, 3, 82; priority of public transport on entry and radial corridors; improvement of railway crossings, bridges and overpasses; landscaping of station and stop approaches; safe pedestrian routes to schools, kindergartens, clinics, shops; continuous cycling and walking connections between residential blocks, parks, and the historic centre; a dedicated programme for poorly accessible peripheral blocks 285, 169, 84, 335, 246.
 
Here we observe a characteristic split: the master plan closes the capital part (network and bridges), FP2MP the service part (public transport, cycling, pedestrian). The two answers do not contradict each other: they target different levels of the transport system and can be compatible under sequential implementation. For example, node~196 in FP2MP, characterised as ``best block by mean accessibility (32.67~min) and maximum connectivity (0.0306)'', logically corresponds to what the master plan already forms in the Aerodrom area (continuation of ul. Novos\"yolov, Bulvar Aviatorov, city-wide arterials).
 
The master plan describes 25 specific facilities in Part~2 ``Parameters of Functional Zones'': 10 reconstructions of MBDOU (kindergartens No.~1, 3, 4, 5, 10, 12, 13, 18, 22), 5 new preschool institutions of 160--220 places, 4 general-education institutions of 600, 825, and 1175 places in V``ezd, Aerodrom (two blocks), Zayachy Remiz and Khokhlovo Pole, 3 school reconstructions (SOSh No.~1, 4, 9, Gatchina Lyceum No.~3), 3 polyclinic facilities (general-practice centre in V``ezd and Roshchinsky, first-type feldsher-obstetric station in Marienburg, dental polyclinic in Aerodrom, 120-bed diagnostic and treatment unit as part of GBUZ LO ``Gatchinskaya KMB'' in V``ezd), 12 sports facilities (including an 8\,732\,m\textsuperscript{2} football indoor arena for 1\,200 visitors, a 5\,000-seat ice arena, an indoor ice rink for 1\,000 spectators, a swimming pool, a tennis court, a 30$\times$18~m physical-fitness complex at ul. Chekhova~9a, reconstruction of the ``Spartak'' stadium of 20\,010\,m\textsuperscript{2}), and 3 cultural facilities (aviation museum in Aerodrom, club-type institution in Centre).
 
FP2MP speaks about different facilities. Its first-wave priority (0--3 years) is to ``close primary care, neighbourhood shops, sports grounds, child services, stops, pedestrian connections'' in blocks 331, 108, 289, 87, 89, 9, 18, 2. FP2MP explicitly mentions none of the 25 facilities of the master plan, because the BlocksNet model does not store references to government construction programmes. However, the majority of master plan facilities are located exactly in the microdistricts that FP2MP treats as priorities: in V``ezd, Aerodrom, Zayachy Remiz, Roshchinsky, and Centre. We record that two-thirds of the deficits appearing in the FP2MP output (clinics, neighbourhood shops, sports grounds) could be closed by master plan facilities, possibly without additional outlays. A simple cross-check---a table ``FP2MP deficit $\to$ master plan facility closing it''---should be the next integration step.
 
The master plan creates an event infrastructure: eleven of the twelve large sports facilities are concentrated in the Aerodrom microdistrict. FP2MP demands elimination of the deficit of sports grounds (provision 0.460) in residential blocks. These are different infrastructure levels: event-related (for spectators and competitions) and everyday (for residents). At this point we see not a disagreement but a division of labour: the master plan solves the event level and does not solve the everyday level, FP2MP sees only the everyday level and does not solve the event level.
 
Both documents agree on the general principle: the historic park core and the palace-and-park ensemble are to be protected. The difference is in the content of the protection. The master plan fixes the constraint through functional zones: 934~ha of public greenery, 149.5~ha of regulated territories, 124.4~ha of water bodies. FP2MP introduces a direct prohibition in its ``Historic--cultural core'' clause---established by NormativeAgent as the protected status of the palace-and-park ensemble as part of a UNESCO object [NormativeAgent sq\_005, SP 42.13330 service radius standards]---and translated into spatial policy: ``do not place large new residential blocks in the influence zone of the palace-and-park ensemble'', ``restrict height and volume where they may disrupt historical scale or visual corridors'', ``preserve the continuity of parks, water bodies, and green links'', ``manage tourist flows without displacing everyday life''.
 
This is substantive agreement on function with a difference in instrument: the master plan protects the core through zoning, FP2MP through direct prohibitions in spatial policy. Both approaches are legitimate. In joint use the master plan provides legally binding frames, FP2MP provides an operational filter that prunes scenarios violating these frames.
 
The master plan and FP2MP use different models of spatial development. The master plan defines ten functional zones with parameters: multi-storey 60.7~ha, mid-rise 135.5~ha, low-rise 59.1~ha, individual housing 428~ha, residential 259.8~ha, multifunctional public-business 122.4~ha, specialised public 48.1~ha, engineering 30.3~ha, transport 258.2~ha, greenery 934~ha, forests 241.3~ha, cemeteries 17~ha, special-purpose greenery 55.4~ha, regulated territories 149.5~ha, water bodies 124.4~ha. This is a monocentric city with a concentrated core and a periphery.
 
FP2MP develops a polycentric network: five territory types, eight priority residential blocks, five service centres (136, 135, 177, 3, 82), one transport-business node (196), five poorly accessible peripheral blocks (285, 169, 84, 335, 246). The historic park core is treated as a separate type, protected from large new volume.
 
This is a conceptual difference, not a technical error. The master plan fixes district planning: where and how many storeys. FP2MP answers a functional question: where to concentrate which functions. Both answers are legitimate in their own coordinates. Alignment requires a cross-check: the facilities of the master plan's functional zones in the 122.4~ha public-business zone should correspond to the service centres 136, 135, 177, 3, 82 of FP2MP. If geometries coincide, the models are aligned. If not, explicit correction is required.
 
\subsubsection{Where the Documents Agree}
\label{sec:comparison-agreements}
 
The aggregate alignment shows partial but not complete correspondence. Where the documents agree in substance, the agreement arises from substantive convergence rather than from a common source:
 
\begin{enumerate}
  \item Protection of the historic park core. Both documents fix a prohibition on large new volume in the influence zone of the palace-and-park ensemble. The master plan enforces this through regulated territories and functional-zone constraints; FP2MP through direct prohibitions in spatial policy (``do not place large new residential blocks in the influence zone\ldots''). Joint use is straightforward: the master plan provides legally binding frames, FP2MP provides the operational filter.
 
  \item Development of social infrastructure in residential microdistricts. The 25 facilities of the master plan are located in V``ezd, Aerodrom, Zayachy Remiz, Roshchinsky, and Centre---the same microdistricts that FP2MP flags as priorities through its \texttt{suggest\_target\_blocks} procedure (blocks 331, 108, 289, 87, 89, 9, 18, 2). The \texttt{compute\_scenario\_provision} pipeline of BlocksNetAgent validates this convergence: for blocks [2, 9, 18, 87, 89, 108, 289, 331], before/after comparison shows clinics improving from 0.295 to 0.833 (13 $\to$ 86 fully served blocks), schools 0.774 $\to$ 0.898, pharmacies 0.821 $\to$ 0.900, sports facilities 0.460 $\to$ 0.545.
 
  \item Node-level consistency. Node~196 in FP2MP (mean accessibility 32.67~min, connectivity 0.0306, derived from the precomputed $351\times351$ accessibility matrix) falls within the Aerodrom area, where the master plan forms a transport corridor (continuation of ul. Novos\"yolov, Bulvar Aviatorov, city-wide arterials).
 
  \item Transport division of labour. The master plan creates the network and bridges, FP2MP creates public transport, cycling, and pedestrian links; these levels complement each other without overlap.
 
  \item Priority-block scale and master-plan microdistrict scale. FP2MP's largest priority block, 331 (11\,087 inhabitants, 505\,928~m\textsuperscript{2} build floor area, 352\,308~m\textsuperscript{2} living area, Shannon diversity 2.7694), sits within the multi-storey zone of 60.7~ha allocated by the master plan, where 12-storey development with 40\% built-up coverage is foreseen. The unit ``priority block'' and the unit ``multi-storey zone'' are compatible: they describe the same territory at different granularity.
\end{enumerate}
 
\subsubsection{Where the Documents Diverge}
\label{sec:comparison-divergences}
 
Where the documents diverge, the divergence is methodological rather than factual:
 
\begin{enumerate}
  \item Demographic basis. The master plan is 112.0 thousand by 2040, FP2MP is 83.3 thousand in the local model. These are two different operational bases. FP2MP operates with 83.3 thousand not because it ignores official data, but because the BlocksNet block model contains exactly this number; the official municipal estimate of 96.3 thousand in 2021 is closer to the master plan's trajectory but is not block-level and therefore cannot be used for spatial prioritisation.
 
  \item Housing expansion. The master plan allocates 683.3~ha of housing. FP2MP, in its ``Options I weighed'' section, explicitly rejects mass peripheral expansion: block~285 with an area of 1\,678\,249\,m\textsuperscript{2} is not developed as a new-housing zone even in the third wave (7--12 years).
 
  \item Spatial-development model. The master plan is a monocentric city with a concentrated core. FP2MP is a polycentric network of subcentres (136, 135, 177, 3, 82) and residential priorities (331, 108, 289, 87, 89, 9, 18, 2).
 
  \item Engineering infrastructure. The master plan contains 83 facilities with parameters (BMK, GRP, water mains, treatment facilities); FP2MP does not include engineering as an independent direction and uses ``pass an engineering check'' as a conditional permit for peripheral development.
 
  \item Everyday services. FP2MP places ``neighbourhood shops, pharmacies, stops, pedestrian links'' at the top of the first wave; the master plan does not contain these facilities as a separate direction.
 
  \item Planning detail. The master plan fixes specific facilities with addresses and capacities; FP2MP works at the level of blocks and provision (e.g.\ \texttt{compute\_service\_provision} returns clinic provision 0.295, with 332 out of 351 blocks underserved; the scenario for blocks [2, 9, 18, 87, 89, 108, 289, 331] increases clinic provision from 0.295 to 0.833).
\end{enumerate}
 
These divergences do not invalidate either document. They reflect different operational models: the master plan normatively fixes the future, FP2MP diagnoses the present.
 